\chapter[HMC Research Program For Nonlinear SSMs]{HMC as a Research Program for Nonlinear State-Space Models}
\label{ch:bf-highdim-hmc-research}

\section{Chapter Claim}
\label{sec:bf-highdim-hmc-claim}

Hamiltonian Monte Carlo for nonlinear state-space models is a per-model
research problem.  A finite scalar target, a dense mass matrix, or a short chain
is not enough.  Nonlinear filters introduce approximate likelihoods, invalid
regions, branch-dependent gradients, and curvature that can defeat plain
Euclidean HMC.  BayesFilter's defensible production direction is an inspectable
fixed-trajectory HMC layer with explicit value/score contracts; TFP NUTS is
diagnostic/reference only in this project.

\section{Target, Coordinates, And Correction}
\label{sec:bf-highdim-hmc-target-gradient}

Let $q\in\mathbb R^d$ be the unconstrained coordinate used by the integrator and
let $\theta=\tau(q)$.  If the filter supplies log likelihood
$\ell_T(\theta)$, the exact potential in $q$ coordinates is
\begin{equation}
\label{eq:bf-highdim-hmc-potential}
  U(q)
  =
  -\ell_T(\tau(q)) - \log \pi(\tau(q))
  - \log\left|\det D\tau(q)\right|.
\end{equation}
If the filter supplies $\widehat\ell_T$, the HMC target is the corresponding
approximate posterior unless an exact correction is supplied.  The gradient
used by the integrator must be the gradient of the same scalar potential used
in the Metropolis correction:
\begin{equation}
\label{eq:bf-highdim-hmc-gradient-contract}
  \nabla_q U(q)
  =
  -D\tau(q)^\top \nabla_\theta
  \{\ell_T(\theta)+\log\pi(\theta)\}\big|_{\theta=\tau(q)}
  - \nabla_q\log|\det D\tau(q)|.
\end{equation}

\begin{proposition}[Value/score same-scalar condition]
\label{prop:bf-highdim-same-scalar}
For a Metropolis-corrected HMC transition to target the density proportional to
$\exp\{-U(q)\}$, the numerical trajectory may use an approximate integrator, but
the acceptance probability must evaluate the same Hamiltonian
$H(q,p)=U(q)+\frac12p^\top M^{-1}p$ whose gradient drives the integrator.
\end{proposition}

\begin{proof}
The standard HMC correction applies Metropolis--Hastings to a deterministic,
volume-preserving, reversible proposal map generated by the numerical
integrator; see \citet{neal2011mcmc} and \citet{betancourt2017conceptual}.  If
the gradient corresponds to a different scalar than the accepted Hamiltonian,
the proposal is no longer the intended discretization of that Hamiltonian and
the usual detailed-balance argument no longer applies.
\end{proof}

\section{Why Plain Mass-Matrix HMC Can Diverge}
\label{sec:bf-highdim-plain-hmc-failure}

A diagonal or dense mass matrix handles global scale but not arbitrary target
geometry.  It cannot by itself remove funnels, hard constraints, invalid
regions, sharp ridges, local curvature rotations, discontinuous branches, or
filter-induced roughness.  Divergences therefore matter: they indicate that the
numerical Hamiltonian trajectory is not faithfully following the target at the
chosen step size and trajectory length.  Speed comparisons are meaningless
until divergence and validity gates pass.

The BayesFilter Model B/C HMC ladder remains a candidate diagnostic only:
finite values and samples were observed in short rows, but acceptance was near
1.0 and maximum $\widehat R$ was around 2.0.  That is evidence that the target
can be evaluated in the tested cells and that research should continue.  It is
not convergence evidence.

\section{Fixed-Trajectory HMC Baseline}
\label{sec:bf-highdim-fixed-hmc}

\begin{algorithm}
\caption{BayesFilter-owned fixed-trajectory HMC diagnostic step}
\label{alg:bf-highdim-fixed-hmc}
\begin{algorithmic}[1]
\Require Potential $U(q)$, gradient $\nabla U(q)$, mass matrix $M$, step size
  $\epsilon$, leapfrog steps $L$, current $q$.
\State Check value and gradient finite at $q$.
\State Draw $p\sim\mathcal N(0,M)$.
\State Run $L$ leapfrog steps with the declared invalid-region policy.
\State Evaluate initial and proposed Hamiltonians with the same scalar $U$.
\State Accept or reject by the Metropolis probability.
\State Record divergence flag, energy error, acceptance, finite status, and
  invalid-region hits.
\end{algorithmic}
\end{algorithm}

This baseline is intentionally simple.  It makes failures visible and keeps
implementation ownership inside BayesFilter.  A production candidate would need
adaptation, diagnostics, and careful warmup, but the research ladder starts with
the smallest transition whose target contract can be audited.

\section{Diagnostic-Only NUTS Policy}
\label{sec:bf-highdim-nuts-policy}

NUTS \citep{hoffman2014nuts} is an important reference algorithm, and Stan's
implementation is the benchmark many practitioners trust.  In this project,
TFP NUTS is not the strategic production path for industrial nonlinear SSMs:
dynamic control flow, adaptation internals, and nested kernel structure make
full-pipeline XLA behavior, failure triage, and model-specific diagnostics less
inspectable than a BayesFilter-owned fixed transition.  This is an engineering
policy, not a mathematical claim that NUTS is inferior.

\section{NeuTra And Transport-Preconditioned HMC}
\label{sec:bf-highdim-neutra}

NeuTra-style HMC learns an invertible neural transport so that HMC runs in a
warped space; local ResearchAssistant identifies \citet{hoffman2019neutra} as
the relevant source but marks the local summary as needing manual review.  The
safe claim is therefore bounded: neural transport is a geometry-preconditioning
idea, not a correctness substitute.

If $q=T_\phi(z)$ and the target density in $q$ coordinates is $\pi_q(q)$, the
warped target in $z$ coordinates is
\begin{equation}
\label{eq:bf-highdim-neutra-target}
  \log \pi_z(z)
  =
  \log \pi_q(T_\phi(z)) + \log|\det DT_\phi(z)|.
\end{equation}
HMC in $z$ must use the gradient of \eqref{eq:bf-highdim-neutra-target} and a
Metropolis correction for that same scalar.  A trained map can reduce curvature
or can make tails worse; only sampler diagnostics decide.

\section{RMHMC, HNN, And Learned Dynamics}
\label{sec:bf-highdim-hmc-acceleration}

Riemannian HMC replaces the fixed mass matrix by a position-dependent metric
\citep{girolami2011riemann}.  The conceptual fit is strong when curvature
rotates over the posterior, but the implementation burden is severe: metric
positive definiteness, derivatives of the metric, implicit solves, and
filtering Hessian stability become part of the target contract.

Hamiltonian neural networks \citep{greydanus2019hamiltonian} and learned HMC
variants can learn proposal dynamics.  The local ResearchAssistant learned-HMC
summary is not domain-specific to nonlinear SSMs, so transfer claims are
forbidden.  Learned dynamics can enter BayesFilter only as a proposal with
target-preserving correction and diagnostics stronger than training loss.

\begin{algorithm}
\caption{Acceleration promotion ladder}
\label{alg:bf-highdim-hmc-accel-ladder}
\begin{algorithmic}[1]
\Require Baseline fixed HMC result, candidate accelerator, controlled target.
\State Verify same-scalar value/score parity in eager and compiled modes.
\State Run plain fixed HMC until validity diagnostics are interpretable.
\State Add transport, RMHMC metric, or learned proposal with exact correction.
\State Reject the acceleration if divergences, nonfinite values, unavailable
  $\widehat R$/ESS, or posterior-reference checks fail.
\State Compare runtime per effective sample only after validity gates pass.
\end{algorithmic}
\end{algorithm}

\section{Diagnostics}
\label{sec:bf-highdim-hmc-diagnostics}

The minimum nonlinear SSM HMC diagnostic ledger is:
\begin{itemize}
  \item finite target and finite gradient at representative points;
  \item value/score parity for the scalar used by the sampler;
  \item divergence count and energy error;
  \item acceptance distribution, step-size sensitivity, and trajectory-length
    sensitivity;
  \item $\widehat R$, bulk/tail ESS, and E-BFMI where chains are long enough for
    these fields to be meaningful;
  \item posterior recovery on a controlled reference before scientific claims;
  \item runtime per effective sample only after the validity vetoes pass.
\end{itemize}
Any divergence, nonfinite value, failed same-scalar check, unavailable
$\widehat R$/ESS on a claimed convergence row, or failed posterior-reference
check reroutes the result to repair.  It must not be ranked by speed.

\section{XLA And GPU Boundary}
\label{sec:bf-highdim-hmc-xla-gpu}

XLA and GPU execution are performance mechanisms, not scientific validation.
For HMC, XLA readiness requires static-enough control flow, stable shapes,
same-scalar value/gradient parity after compilation, deterministic invalid
region policy, and reproducible diagnostics.  GPU speedup requires trusted GPU
execution, meaningful problem size, CPU comparator rows, and validity gates
before speed ranking.  CPU-only artifacts must hide GPU devices before
framework import and say so.

\section{Industrial Practitioner Stress Test}
\label{sec:bf-highdim-hmc-industrial}

At NAWM-like scale, plain HMC can fail because posterior geometry is shaped by
structural solution maps, occasionally stiff measurement equations, parameter
constraints, and filter approximation error.  The plausible rescue sequence is:
\begin{enumerate}[label=(\arabic*)]
  \item stabilize the filtering target and value/score path first;
  \item establish a plain fixed-HMC baseline and record its divergences;
  \item use block or transport preconditioning to reduce geometry pathologies;
  \item consider HNN or learned proposals only after exact correction and
    sampler diagnostics are routine;
  \item treat every model as a new research target until posterior checks pass.
\end{enumerate}

\section{Source And Evidence Ledger}
\label{sec:bf-highdim-hmc-ledger}

\begin{longtable}{p{0.24\linewidth}p{0.28\linewidth}p{0.38\linewidth}}
\toprule
Claim & Support class & Boundary \\
\midrule
\endhead
HMC same-scalar and diagnostic discipline & Primary sources
\citet{neal2011mcmc,betancourt2017conceptual} plus Proposition~\ref{prop:bf-highdim-same-scalar}
& Does not validate any BayesFilter nonlinear chain. \\
NUTS is diagnostic/reference-only here & Project policy informed by
\citet{hoffman2014nuts} and local engineering constraints & Not a claim against
Stan or NUTS mathematically. \\
NeuTra is geometry preconditioning & \citet{hoffman2019neutra}; local
ResearchAssistant summary is needs-review & No BayesFilter NeuTra validation or
Gate-1 promotion in this lane. \\
RMHMC is geometry-aware & \citet{girolami2011riemann}; local summary
needs-review & Hessian/metric implementation unresolved. \\
HNN/learned dynamics are proposal acceleration ideas & \citet{greydanus2019hamiltonian}
and local learned-HMC summary & No nonlinear SSM transfer claim. \\
Model B/C HMC ladder is not convergence evidence & BayesFilter V1 result notes
& Finite short rows nominate research only. \\
\bottomrule
\end{longtable}

\section{Unresolved Claim Register}
\label{sec:bf-highdim-hmc-unresolved}

\begin{itemize}
  \item No BayesFilter nonlinear HMC convergence result exists.
  \item No NeuTra, RMHMC, or HNN BayesFilter backend is implemented or promoted
    here.
  \item No nonlinear Hessian readiness or NAWM HMC diagnostic is claimed.
  \item ResearchAssistant summaries for NeuTra/RMHMC/learned HMC are useful
    review material but remain needs-review, so they cannot support
    theorem-level prose.
\end{itemize}

\section{What Is Not Concluded}
\label{sec:bf-highdim-hmc-nonclaims}

This chapter does not conclude that HMC, NeuTra, HNN, RMHMC, TFP NUTS, XLA, or
GPU execution is production-ready for nonlinear SSMs.  It defines the
per-model research ladder required before such a claim could be considered.
